\documentclass[11pt]{article}
\usepackage{amsmath,amssymb}
\usepackage[margin=1in]{geometry}
\usepackage{hyperref}

\title{Ideas Parked --- Not Yet Merged (v6)}
\author{Staging document; each entry marks its eventual destination}
\date{}

\begin{document}
\maketitle

\noindent This document holds ideas discussed against Pattern Arithmetic
v15\_19, Inference, and Assembly, none of which have been written into
any of those three files yet. Each entry states its status and its
intended destination. Nothing here is settled; the point of this file
is to let that be true without leaving Paper 1 mid-edit.
v5 added one scribble: an identity handle on an atom as a meta field,
for an honest ``uncollapse'' that looks up a kept bowl rather than
inverting Mix. v6 adds two: a sphere-packing bound relevant to the
path-reduction item's distance requirement, and a new theoretical
problem --- merging every atom in a unit space to one representative
--- which splits the same way the $\tau_{\text{weather}}$ correction
already did, and is recorded as a second instance of that same pattern
rather than a new result.

\section{Destined for Pattern Arithmetic (still formally open ---
but a leading resolution candidate has emerged, with a citation)}

\paragraph{Mix across differing directions.} The paragraph's own text
still reads ``left open.'' Nothing here changes that text. What
changed is that the paper already contains a worked, concrete instance
of this exact question, in a section not previously connected to it.

\textbf{[Found: the paper already answers this question once, for one
sort, and points toward a general answer.]} The $\mathtt{chladni}$
inclusion states, word for word: ``Across modes the map is not
available, and not for want of effort. Mode space is not closed under
addition... Cross-mode coexistence is therefore $\oplus$ --- a bowl of
modes --- and the resulting figure is a reading of that bowl which is
not $\boxplus$. This is the open question of
[the differing-directions open item] met in a concrete domain,
and it suggests the question may be closed negatively for this sort
rather than left waiting.'' This is a citation, not a new argument ---
the paper's own words, in a section that was not previously read as
bearing on Gap~1 directly.

\textbf{[Why this converges with everything else worked through in
this project, not just chladni.]} Once Mix is read as ``resolve
several reports on one shared question into a verdict'' (argued in
v2/v3, carried forward below unchanged), refusing across-direction
combination stops looking like a gap and starts looking like the
correct behavior of an operator that was never asked that question in
the first place. And every practical workaround that actually worked
in this project was already living on the $\oplus$ side of precisely
this line, independently: bowls holding separate weather channels,
per-channel similarity scoring, the candlestick summary for
reference-atom building, Assembly's own camp-to-system Hill bijection.
None of them ever needed Mix to cross a direction. All four arrived,
separately, at ``keep it in a bowl, read it there'' --- which the
chladni section had already stated as a formal result.

\textbf{[Status: a leading candidate, not yet a decision.]} This
supports Route 1 (refuse outright, from the three already on file)
over Route 2 (rotate first) or Route 3 (two-stage operator) as the
likely resolution, with a real citation behind it rather than
preference. Writing this into Paper 1 would mean editing the
``Open: Mix across differing directions'' paragraph to state the
refusal explicitly, citing the chladni section as precedent. Not done
here, on request --- parked, ready, pending a decision to merge.

\textbf{[Carried forward unchanged: the reading argued in v2/v3.]} Mix
is not ``multiple vectors becoming one vector.'' It takes several
reports already agreed to answer one shared question and returns a
verdict about that question. The shared $\hat n$ is the precondition
that makes ``these answer the same thing'' true, not a boundary
condition on a more general operator. Supporting evidence: same-
direction opposite-phase cancellation to black (principled
contradiction-handling, not a fabricated compromise); unmixing's
proven ill-posedness, with metamerism as independent real-world
precedent; the paper's own opening line, ``two notes yield a chord,
not a larger integer... the word `add' names a different rule.''

\textbf{[Carried forward unchanged: the ``we mix all the time''
challenge, resolved.]} Sugar in tea is one dial, never needed
differing directions. Red plus green additively sits ``at the same
place'' in the paper's own words --- already inside Mix's settled
domain. Paint-layer flattening is confirmed, across every source
checked, as destructive and irreversible, matching unmixing exactly;
compositing unflattened layers is $\oplus$, not $\boxplus$, which is
why ``unmerge them later'' was always possible. Representing a whole
painting as one vector remains Gap 1 itself, not a counterexample to
it.

\textbf{[Carried forward unchanged: the two-part split, now read as
one part likely answered.]} Embedding validity --- does geometric
closeness on a shared sphere track real semantic agreement at all ---
was the harder, prior question, true for hue by deliberate
construction, unstated for an arbitrary sphere of concepts. Given the
chladni precedent above, the likely answer for sorts where this was
never stated is that no valid embedding needs constructing, because
cross-direction combination is refused before the embedding question
would even matter.

\textbf{[Carried forward unchanged: still confirmed not to help.]}
Magnitude self-bounding near opposition does not license the
direction that comes with it. The Hill form and Assembly's Hill
bijection both fix magnitude only, verified twice, never direction.

\section{A correction, recorded so it is not repeated}

\paragraph{$\tau_{\text{weather}}$ names two different problems, and
only one of them is Gap~1.} Raised once as a proposed admissible sort
for combining different weather concepts (rain, drought) across
directions --- that question is genuinely Gap~1, and the discussion
above bears on it. Raised a second time, later, meaning something
narrower: given several already closely-aligned vectors within one
bounded region, how to arrive at a single representative atom, using
the Hill form. That second question does not invoke $\boxplus$ across
differing directions at all --- it is Assembly's already-verified
system-promotion mechanism (16 camps to 13 systems under the Hill
bijection), recognized in a new setting, not a Paper~1 admissibility
question. The two were briefly conflated in this project; recorded
here so the same conflation is not repeated. Assembly's version needs
nothing further; Gap~1's version is addressed in Section~1 above.

\section{Destined for Pattern Arithmetic (advances an existing open item)}

\paragraph{Location as a sort --- a third candidate.} The existing
paragraph names two candidates and adopts neither. A third was tested
here: dedicate separate, disjoint dimensions to content, to time, and
to location, rather than letting any of them share an axis.

\textbf{[Verified.]} Two shifts acting through disjoint content
dimensions commute exactly (checked numerically, not approximately:
time-then-location and location-then-time landed on the identical
point to full floating-point precision). Two shifts forced to share a
single content dimension do not commute, confirming the disjointness
is what matters, not rotation itself.

\textbf{[Verified.]} Content stays recoverable after both shifts:
projecting back onto the content dimensions alone reproduced the
original direction closely, while the time-axis and location-axis
each carried a clean, separately-readable amount of shift, with
nothing blended between them.

\textbf{[Open, unchanged.]} The encoding itself --- what specific
rotation a given time value or location corresponds to --- is not
stated anywhere. The geometry now supports it cleanly; the mapping
does not exist yet.

This is real, checkable content, not yet written as an addition to
the existing paragraph.

\section{Destined for Pattern Arithmetic (open item stays open, unchanged)}

\paragraph{Reducing a path by merging close actors.} The existing
paragraph asks for three things: a stated distance, a merge rule, and
a bound on how far a reduced path's endpoint may drift from the
unreduced one. Discussion here supplied a distance (see below) and
partially addressed the merge rule, but only for magnitude, not for
direction --- and direction is the harder half this item is actually
about.

\textbf{[Verified.]} Switching the strength calculation from the
companion paper's clip to Inference's already-established Hill form,
$r'=|z|/(K+|z|)$, changes nothing about the resulting direction
$\chi'=\arg(z)$. Checked directly on two constructed sequences with
identical content and opposite order: both gave the same angle under
either formula. The Hill substitution fixes evidence-count blindness
in strength; it does not touch the order-blindness this open item is
actually asking about.

\textbf{[New: a candidate grounding for the stated distance,
item~(1) of the three things required.]} How many actors can be
mutually distinguishable --- separated by at least some angle
$\theta_{\min}$ --- on $\mathcal{S}$ at all is a real, named problem
(spherical codes, the Tammes problem), not an arbitrary choice to
make. Exact optimal counts are known only for small $N$; a standard
approximate bound, treating each point as owning a non-overlapping
cap of half-angle $\theta_{\min}/2$, gives
\[
N_{\max} \;\approx\; \frac{4\pi}{2\pi\bigl(1-\cos(\theta_{\min}/2)\bigr)}
\;=\; \frac{2}{1-\cos(\theta_{\min}/2)}.
\]
Computed for illustration:
\begin{center}
\begin{tabular}{rr}
$\theta_{\min}$ & approx.\ max.\ distinct points \\
\hline
$1°$   & $52{,}525$ \\
$5°$   & $2{,}101$ \\
$10°$  & $526$ \\
$30°$  & $59$ \\
$60°$  & $15$ \\
$90°$  & $7$ \\
$120°$ & $4$ \\
\end{tabular}
\end{center}
This is an approximation, not an exact extremal result, but it is a
real, checkable mathematical object rather than a guess: naming a
$\theta_{\min}$ for ``close enough to merge'' directly bounds how many
genuinely distinct actors can exist near each other at all, which
bears on item~(3), the drift bound, as well as item~(1).

\textbf{[Open, unchanged.]} No merge rule for direction exists. No
drift bound exists for either magnitude or direction. This item
remains exactly as open as it was before this discussion, now with
the magnitude half more precisely separated from the direction half,
and a real candidate supplied for the distance half specifically.

\section{Destined for Inference (new mechanism)}

\paragraph{Candlestick summary: a richer merge rule for reference-atom
building.} Not a resolution of the Paper 1 item above --- a different
question. Paper 1's open item is about \emph{walking}: can a shortened
path still be walked to roughly the same endpoint. This is about
\emph{reading}: how to summarize a group of atoms for comparison
against a reference atom, which is Inference's own established
business (its reference-atom section), not $T$'s.

\textbf{[Verified.]} Open/close/high/low, computed the way a trading
candle is, distinguishes exactly the case flat Mix could not. Two
sequences, identical content, opposite order --- naive Mix gave the
identical angle ($47.5°$) for both. The candle's open and close swapped
between them ($20°\!\to\!160°$ versus $160°\!\to\!20°$), correctly
reading one as ending badly and the other as ending well. High and low
stayed identical for both, correctly capturing that only the
\emph{range} was shared, not the \emph{direction of travel}.

\textbf{[Settled, matches existing machinery.]} The scope rule ---
every atom belongs to exactly one time unit, no overlap between
buckets --- is not new; it is precisely the discrete-time-bucket
metadata already defined on a walk's own steps in the companion paper.

\textbf{[Open.]} This gives Inference a genuinely better reference-atom
merge rule than flat Mix. It is not proposed as a fix for Paper 1's
walking-specific open item, which remains untouched above. Whether
open/close/high/low itself needs a coherence check analogous to
$\bar R_c$ has not been considered.

\section{Destined for Inference (naming caution, not yet a mechanism)}

\paragraph{``Inference Mandelbrot'' is not Assembly's Mandelbrot, and
should not share its name.} Assembly's multi-zoom promotes an
unordered thing (a settled camp) and is fully handled by the Hill
bijection already verified there. Whatever would promote an
\emph{ordered} thing (a walked path, a movie, a time series) into one
higher-level atom is a different, harder problem --- it is the
direction half of the path-reduction item above, restated from a
different angle. \textbf{[Naming flag only.]} If this is ever built,
it wants its own name; reusing ``Mandelbrot'' for it would be the same
kind of collision already caught and corrected for ``collapse,''
``singularity,'' and the quantum swap test.

\section{Destined across three papers (new, split by piece)}

\paragraph{Meta vectors: a second kind of vector alongside information
vectors.} Proposed: an information atom may carry two distinct kinds
of vector. Information vectors are Axiom 7's bundle $v_i$'s, capped at
$r_i\le1$, potentially several sorts $\tau$ across a whole record.
Meta vectors describe the atom rather than being part of its content
--- three named so far: $V_{\text{derived}}$, link vectors (pointers
between atoms, one- or two-directional), and a timestamp vector.
Explicitly left open whether more meta-vector types are needed. Meta
vectors are not capped at 1, and grouping several atoms' meta vectors
together on like terms is proposed as a distinct mechanism for
building higher structures. This is one idea that lands in three
different places once traced to where each piece actually belongs.

\textbf{[Destined for Pattern Arithmetic.]} The taxonomy itself ---
information vectors versus meta vectors as a formal, named distinction
sitting on top of Axiom 7 --- is structural, the same kind of addition
Axiom 7 itself was. $V_{\text{derived}}$ already lives there. A
timestamp vector, if adopted, would be a third, genuinely different
treatment of time from the two already on file above: distinct from
the discrete-bucket metadata already defined on a walk's own step, and
distinct from the dedicated-axis encoding in the location-as-a-sort
entry. Not yet clear whether more than one of the three should
coexist, or which.

\textbf{[Destined for Assembly.]} Link vectors specifically. This is
the mechanism the shelf's own neighbour-graph note was already naming
without a concrete form: ``a data structure on seats after rest, not
$T$.'' Grouping link vectors on like terms is collecting edges into a
graph, not merging them --- gathering, in the $\oplus$ sense --- so
it carries none of the open questions below it.

\textbf{[Destined for Inference.]} The grouping mechanism itself:
querying several atoms' meta vectors, confirming they sit close
together, and promoting one shared meta-vector up to the group as a
whole. This is the space-time gating idea from earlier discussion,
generalized into a standing principle. Timestamp and location get
\emph{used} here; they are only \emph{defined} in Pattern Arithmetic
above.

\textbf{[Verified, carried forward as a caveat rather than a new
finding.]} Combining several atoms' $V_{\text{derived}}$'s stays safe
only under the constraint the original already carries: auxiliary,
never promoted to a trusted, walkable atom on its own. Checked
directly that neither Inference's Hill form nor Assembly's Hill
bijection touches this constraint --- both fix magnitude only, never
direction --- so grouping $V_{\text{derived}}$'s inherits that same
limit rather than escaping it.

\textbf{[Note added here.]} Link vectors were separately discussed as
one possible mechanism for a limited, honest form of ``uncollapse'':
not inverting Mix (still impossible, unaffected by anything here),
but keeping an explicit, external pointer from a collapsed atom back
to the bowl that produced it, created at collapse time, before the
bowl itself might otherwise be discarded. This is navigation via a
kept record, not reconstruction from the collapsed result --- the
distinction the white-light/prism discussion turned on, confirmed
there to matter.

\section{Destined across three papers (scribble, v5)}

\paragraph{Identity handle as a meta field, for honest uncollapse.}
An information atom may carry one more meta quantity: an opaque
identity $id$, assigned before any gather or Mix, and kept on the
atom afterwards. The scribble that prompted this note called it a
``meta vector'' with $\tau_{id}$ equal to that identity. Fine-tuning
below refuses $\tau_{id}$ as a sort and keeps $id$ as a \emph{field},
in the same shelf as timestamp, mint-location, and link pointers
already named in the meta-vector entry.

\textbf{[What it is for.]} Mix remains irreversible (Paper 1 v1.1).
``Uncollapse'' here does not reconstruct ingredients from a yellow.
It means: at collapse time, write a pointer from the resulting atom
(or from the register occupant after a walk) back to the bowl that
produced it, and keep that bowl under the same $id$ (or a pair
$result\_id \leftrightarrow bowl\_id$). Later, given the result, look
the bowl up. Navigation by a kept record, not inversion. This is the
same distinction already recorded under link vectors and the
white-light / prism discussion.

\textbf{[Viability --- what works.]}
An opaque handle is the right shape. It does not need to live on
$\mathcal{S}$, does not need $r\le 1$, and does not need a direction.
A string or integer key is enough. Creating the pointer \emph{before}
the bowl is discarded is the whole trick; after discard there is
nothing to uncollapse. Links already named that graph; $id$ is the
node name those edges need. Assembly can own the edge table
(neighbour / parent-bowl). Paper 1 only needs the taxonomy: $id$ is
meta, not a unit of truth. Inference may use $id$ to re-open a bowl
and run a second reading.

\textbf{[Viability --- what to refuse now.]}
Do not mint $\tau_{id}$. $\tau$ is the legality tag of an information
vector. An identity is not a pattern and must not appear in the type
table. Do not encode $id$ as an angle on the sphere; colliding two
ids geometrically would be a false Mix. Do not claim that
$\{A\}\ominus A$ or $T^{-1}$ plus $id$ inverts $\boxplus$. Gather
still satisfies $A\oplus A=A$: two tokens of the same pin share one
address and, if they are the same \emph{instance}, one $id$; two
distinct instances of the same pin (two physical chairs) need two
ids and are not collapsed by $\oplus$ unless the lexicon says they
are the same state. That instance-versus-kind cut is the fine-tuning
this scribble does not finish.

\textbf{[Tiny picture.]}
Bowl $B=\{\mathrm{red},\mathrm{green}\}$, $id_B=\#17$.
Mix reads yellow, $id_Y=\#18$, meta link $\#18\to\#17$.
Uncollapse($\#18$) returns $B$, still red and green, still two atoms.
The yellow is not pulled apart. The prism was kept on the shelf.

\textbf{[Status.]} Scribble. Destined: taxonomy to Paper 1 (meta
shelf, next to $V_{\mathrm{derived}}$); the pointer table to
Assembly; lookup-and-reread to Inference. Not merged. Fine-tune
later: who assigns ids, instance vs kind under $A\oplus A=A$, and
whether one atom may keep a stack of parent-bowl ids after several
collapses.

\section{A theoretical problem, split the same way Section~2 already
split one --- new, v6}

\paragraph{Can every atom in a unit space, across all labels and
types, be merged to one representative atom?} Posed precisely: a
bounded region holds many already-settled atoms, of possibly
\emph{differing} sorts $\tau$, not just differing content within one
sort. Can the whole collection become one representative for that
region? Framed, as requested, with both a theoretical and a practical
answer in view, the way the sphere-packing question above was.

\textbf{[Solved --- but only for the restricted case, which is
Section~2's correction in general form.]} When every atom in the
region shares one sort (or, more precisely, one already-admitted
$\varphi_\tau$ direction-family), this is exactly
Assembly's already-verified system-promotion mechanism: many close
atoms, one representative, via the Hill bijection rather than the
naive clip, because Hill preserves evidence-count where the clip
saturates and loses it. Sixteen camps correctly became thirteen
systems on this basis, not eleven. ``A lot of random atoms'' changes
nothing here --- randomness in position or content is fine, provided
the sort is shared. This is Section~2's $\tau_{\text{weather}}$
correction, restated as the general pattern it always was: random
instances of one $\tau$, closely aligned, merge cleanly.

\textbf{[Still open --- this is Gap~1, not a second result.]} ``Across
all labels and types,'' taken literally, names the harder case: atoms
of genuinely \emph{different} sorts in the same region. This is not
solved by Hill's equation or the bijection. Confirmed twice already
in this project: both Inference's Hill form and Assembly's Hill
bijection fix magnitude saturation only, and neither touches direction
or sort. Merging across genuinely different sorts into one atom is
precisely Gap~1 (Section~1), for which a leading candidate (refuse
outright, per the $\mathtt{chladni}$ precedent) exists but has not
been written into Paper~1.

\textbf{[Theoretical versus practical, for the cross-type half
specifically.]} The sphere-packing question above had a well-defined
continuous quantity (point count) that a practical approximation could
meaningfully estimate. This question's open half does not: there is no
stated $\varphi_\tau$ for combining across sorts to approximate in the
first place. A random or approximate merge rule cannot stand in for a
merge rule that has not been defined --- this is an admissibility gap,
not a precision gap, and the two should not be read as the same kind
of ``open'' question.

\textbf{[So, precisely, correcting the framing in which this was
posed.]} ``We solved it using bijection and Hill's equation'' is true
for same-sort regions, and is a real, verified, reusable result ---
but it is not yet true for atoms spanning genuinely different types,
which remains exactly Gap~1, unresolved by anything checked so far.
The two cases should be kept as separate claims, the same way
Section~2 already separated them once for $\tau_{\text{weather}}$
specifically; this entry generalizes that same correction rather than
adding a new one.

\section{Considered and found to need nothing filed}

For completeness: the reference-atom and precursor-witness mechanism,
the bundle-vector / uncapped $V_{\text{derived}}$ distinction from
Mix, the $\pi=(x;w_1,\ldots,w_n)$ poem trajectory, and the precise
definition of $\chi$ as internal phase-color (already stated in
Axiom 1: ``two identities may occupy the same direction without
overwriting one another'') were all discussed at length and all
turned out to already be fully present in Inference and Paper 1
respectively. Nothing from those threads is parked here, since there
was nothing left unresolved to park.

\end{document}
